\chapter{Configuration Space}
\label{ch:configuration-space}

The configuration space is the unifying abstraction of robot motion:
once it is in place, planning, control, and localization all become
statements about the same underlying object. This chapter introduces
the concept, its structure, and its role in motion planning.

\section{Motivation: Why Configuration Space?}
\label{sec:why-cspace}

The central problem of mobile robotics is the following: given a robot
in an environment with obstacles, move it from an initial configuration
to a final one without collisions. Stated directly in the workspace,
this problem is hard: the robot has an extent, the obstacles have
arbitrary shapes, and checking whether a given placement is collision-free
requires geometric reasoning.

The \emph{configuration space} (C-space), introduced by Lozano-Pérez in
the early 1980s, makes the problem tractable. The idea is to replace
the motion of an extended body among obstacles with the motion of a
\emph{point} in an abstract space, where obstacles have been ``inflated''
to account for the robot's shape. This reduction turns a complex
geometric problem into a path-search problem in a topological space,
for which a large body of algorithms exists
\cite{choset2005principles, siciliano2025foundations}.

\section{Configuration and Degrees of Freedom}
\label{sec:configuration-dof}

\begin{definition}[Configuration]
A \emph{configuration} \( q \) is a complete specification of the
position of every point of the robot. It is represented as a vector
\( q = [q_1, \dots, q_n]^\top \), where \( n \) is the number of
\emph{degrees of freedom} (DOF) of the robot.
\end{definition}

The degrees of freedom are the minimum number of independent parameters
needed to determine the position of the robot in space. They are a
property of the robot's kinematics, not of the environment.

\section{The Configuration Space}
\label{sec:cspace-definition}

\begin{definition}[Configuration Space]
The \emph{configuration space} \( C \) is the set of all possible
configurations of the robot. Formally, \( C \subseteq \mathbb{R}^n \),
although in many cases it carries a richer structure: for instance,
\( C = \mathbb{R}^2 \times \mathcal{S}^1 \) for a planar robot with
orientation, where \( \mathcal{S}^1 \) denotes the unit circle.
\end{definition}

The dimension of \( C \) equals the number of DOF of the robot, which
in general is \emph{not} the dimension of the workspace. This is a
common source of confusion at first.

\section{Workspace vs Configuration Space}
\label{sec:workspace-vs-cspace}

It is essential to distinguish two spaces:

\begin{itemize}
    \item The \textbf{workspace} \( \mathcal{W} \) is the physical space
          in which the robot moves (e.g.\ \( \mathbb{R}^2 \) or
          \( \mathbb{R}^3 \)), expressed in world coordinates.
    \item The \textbf{configuration space} \( C \) is the abstract space
          of all configurations the robot can assume. Its dimension
          equals the robot's DOF.
\end{itemize}

A planar robot that can both translate and rotate has a 3-dimensional
C-space, even though it lives in a 2-dimensional workspace. A free
rigid body in \( \mathbb{R}^3 \) has a 6-dimensional C-space, while
its workspace is 3-dimensional.

\section{Examples}
\label{sec:cspace-examples}

\subsection{Point Robot}
\label{subsec:point-robot}

If the robot is a point translating in \( \mathbb{R}^2 \), then
\( q = (x, y) \) and \( C = \mathbb{R}^2 \). The C-space coincides
with the workspace, and the C-obstacles are the original obstacles.

\subsection{Disk Robot}
\label{subsec:disk-robot}

A disk-shaped robot of radius \( r \) translating in \( \mathbb{R}^2 \)
has \( q = (x, y) \), where \( (x, y) \) is the center of the disk.
Again \( C = \mathbb{R}^2 \), but each obstacle is \emph{inflated} by
the robot's radius: a point obstacle becomes a disk of radius \( r \),
and a polygonal obstacle is expanded by the \emph{Minkowski sum} with
the robot's disk. After this inflation, the disk robot can be treated
as a point in \( C \).

\subsection{Robot with Orientation: \( \mathbb{R}^2 \times \mathcal{S}^1 \)}
\label{subsec:cspace-orientation}

If the robot is a disk that can also rotate, then
\( q = (x, y, \theta) \) and \( C = \mathbb{R}^2 \times \mathcal{S}^1 \),
where \( \mathcal{S}^1 \) is the unit circle (the angle \( \theta \) is
periodic). This space has a non-trivial topology: it is \emph{not}
\( \mathbb{R}^3 \), because \( \theta \) and \( \theta + 2\pi \) denote
the same configuration. The intuitive picture is a \emph{cylinder}:
one linear dimension \( (x, y) \) and one circular dimension
\( (\theta) \).

\subsection{Planar 2-Link Manipulator: \( \mathcal{T}^2 \)}
\label{subsec:cspace-2link}

For a planar manipulator with two revolute joints,
\( q = (\theta_1, \theta_2) \) and
\( C = \mathcal{S}^1 \times \mathcal{S}^1 = \mathcal{T}^2 \), the
\emph{2-dimensional torus}. Each joint is a circle, and the product
of two circles is a torus. This is a classic example of how the
topology of \( C \) can differ radically from that of
\( \mathbb{R}^n \).

\subsection{Rigid Body in 3D: \( SE(3) \)}
\label{subsec:cspace-se3}

A free rigid body in \( \mathbb{R}^3 \) has 6 DOF: three translational
and three rotational. Its C-space is
\( C = \mathbb{R}^3 \times SO(3) = SE(3) \), the \emph{special
Euclidean group} of dimension 3. The rotational part \( SO(3) \) is
not \( \mathbb{R}^3 \): it is a compact 3-dimensional manifold.

\section{C-Obstacles and Free Space}
\label{sec:c-obstacles}

Once \( C \) is defined, obstacles from the workspace must be mapped
into it.

\begin{definition}[C-Obstacle]
Given an obstacle \( O_i \subset \mathcal{W} \), the
\emph{C-obstacle} \( CO_i \) is the set of configurations in which
the robot collides with \( O_i \):
\[
CO_i = \{ q \in C : B(q) \cap O_i \neq \emptyset \},
\]
where \( B(q) \) is the volume occupied by the robot at configuration
\( q \).
\end{definition}

The union of all C-obstacles is the \emph{C-obstacle region}:
\[
CO = \bigcup_i CO_i.
\]

\begin{definition}[Free Configuration Space]
The \emph{free configuration space} is the complement of \( CO \) in
\( C \):
\[
C_{\text{free}} = C \setminus CO
= \{ q \in C : B(q) \cap \textstyle\bigcup_i O_i = \emptyset \}.
\]
\end{definition}

Planning a safe motion is equivalent to finding a path in
\( C_{\text{free}} \) connecting the start configuration \( q_s \) to
the goal configuration \( q_g \).

\paragraph{The point-robot trick.}
If the robot is modeled as a point, the workspace obstacles must be
inflated by the Minkowski sum with the robot's geometry. After this
transformation, planning for the extended robot reduces to planning
for a point in \( C_{\text{free}} \). This is the canonical formulation
of the motion planning problem.

\paragraph{Path vs.\ trajectory.}
A \emph{path} is a purely geometric sequence of configurations, with
no timing information. A \emph{trajectory} is a path endowed with a
time law \( q(t) \). In this chapter, and in motion planning in
general, the object of interest is the path. Timing becomes relevant
only when the dynamics and the control layer are brought into play.

\section{Relevance to Motion Planning}
\label{sec:cspace-motion-planning}

The configuration space reformulates the motion planning problem in
a clean and general way:

\begin{quote}
Given \( C \), a start configuration \( q_s \in C_{\text{free}} \),
and a goal configuration \( q_g \in C_{\text{free}} \), find a
continuous path \( \tau : [0, 1] \to C_{\text{free}} \) such that
\( \tau(0) = q_s \) and \( \tau(1) = q_g \).
\end{quote}

All the algorithms studied later in the course --- roadmaps, cell
decompositions, probabilistic planners, potential fields --- are
methods for solving this problem in \( C_{\text{free}} \). The
differences between them lie in how they represent \( C_{\text{free}} \)
and how they search it.

A subtle point: the \emph{topology} of \( C \) matters. Searching for
a path in \( \mathbb{R}^n \) is different from searching in
\( \mathcal{S}^1 \), \( \mathcal{T}^n \), or \( SE(3) \). Planners
must respect the topology of the space in which they operate.

\section{Takeaways}
\label{sec:cspace-takeaways}

\begin{enumerate}
    \item A \textbf{configuration} \( q \) is a complete specification
          of the robot's position.
    \item The \textbf{configuration space} \( C \) is the set of all
          configurations; its dimension equals the robot's DOF.
    \item The \textbf{workspace} is where the robot lives; the
          \textbf{C-space} is where planning happens.
    \item \textbf{C-obstacles} are the image of workspace obstacles
          in \( C \); \textbf{\( C_{\text{free}} \)} is the set of
          collision-free configurations.
    \item Motion planning reduces to finding a path in
          \( C_{\text{free}} \) from \( q_s \) to \( q_g \).
    \item The \textbf{topology} of \( C \) matters: \( \mathbb{R}^n \),
          \( \mathcal{S}^1 \), \( \mathcal{T}^n \), \( SE(3) \) are all
          different spaces, and planners must respect their structure.
\end{enumerate}

\section{References}
\label{sec:cspace-references}

The material in this chapter follows
\cite{siciliano2025foundations} for the foundational definitions and
\cite{choset2005principles} for the algorithmic treatment of
C-obstacles and free space.